%!TEX root = ../main.tex
\section{Conclusion}

This paper presented the architectural design and implementation of a workflow-enabled enterprise document and records management system that integrates compliance-oriented modeling directly into its core domain structure. 
Rather than treating workflow orchestration and records management as independent subsystems, the proposed solution unifies classification hierarchies, lifecycle attributes, configurable metadata, and versioned workflow execution within a single coherent architecture.

The presented model demonstrates that regulatory alignment, structural consistency, and process flexibility can coexist without introducing architectural fragmentation. 
By separating workflow definitions from runtime execution state, enforcing schema-driven metadata validation, and embedding retention logic into classification entities, the system establishes a compliance-first foundation while preserving configurability and extensibility.

The architectural approach further illustrates how modern service-oriented design, asynchronous processing, and container-based deployment can support scalable and maintainable enterprise registry solutions. 
The hybrid metadata persistence strategy and versioned workflow model contribute to both adaptability and auditability, which are essential characteristics of regulated information systems.

Future work will focus on automated lifecycle enforcement mechanisms, extended compliance validation layers, and empirical performance evaluation under large-scale distributed workloads. 
Additional research may also explore generalization of the proposed model to other regulated domains requiring traceable document-centric workflows.
